% authority: science_draft
% claim_status: draft/control
\documentclass[11pt]{article}
\usepackage[margin=1in]{geometry}
\usepackage{amsmath,amssymb,booktabs,array,enumitem,hyperref}
\hypersetup{hidelinks}
\setlist{nosep}

\title{V22 P1-T03 Matter Position and No-Postulate Gate Policy\\
\large Position A conditional input, Position B reopening, and evidence-only Gates B/D}
\author{Smuggling Auditor parent with Ontology Formalizer perspective}
\date{2026-08-09}

\newcommand{\Gate}{\operatorname{Gate}}
\newcommand{\Supp}{\operatorname{Supp}}
\newcommand{\Positive}{\mathsf{POSITIVE}}
\newcommand{\NotReady}{\mathsf{NOT\_READY}}
\newcommand{\Insufficient}{\mathsf{CANDIDATE\_INSUFFICIENT}}

\begin{document}
\maketitle

\begin{center}
\fbox{\parbox{0.92\textwidth}{
\textbf{Status and authority.} This is a draft/control integration decision for
V22 P1-T03. It preserves the exact protected P7 finite source-matter adoption
without extending it. It is not a new protected Gate verdict, ontology edit,
matter derivation, target coupling, Einstein-equation derivation, benchmark
promotion, external review, or completed derivation.
}}
\end{center}

\section{Decision: Position A for the next geometry cycle}

Position A is adopted as the planning boundary for the next geometry cycle:
\emph{a fixed matter theory is a declared conditional input to geometry
reconstruction}. The only current matter input authorized by this decision is
the exact, unchanged P7-T01 through P7-T06 finite package under the protected
P7-T08 decision and its eight constitutive postulates. ``Input'' means that a
later, separately admitted task may assume those exact source-side structures
inside their declared finite domains. It does not mean that realistic matter,
geometry, universal target coupling, or gravitational dynamics has been
derived.

The decision is non-retroactive and non-expansive. It does not alter any P7
formula, hash, finite domain, interpretation label, obstruction, or forbidden
overread. The open obstruction
\texttt{OBST-P7T07-CROSS-LAYER-COMPOSITION-GAP-001} remains an
\texttt{open\_derivational\_gap\_after\_constitutive\_adoption}.

\subsection{Exact protected package snapshot}

\begin{center}
\small
\begin{tabular}{@{}p{0.26\textwidth}p{0.56\textwidth}@{}}
\toprule
Registered object & Exact SHA-256 \\
\midrule
P7-T01 source-matter ontology & \path{8d160217bf223078a11bc63fde6593c11c39d5b50d9c48fbad7b12084f8a752d} \\
P7-T02 finite transition kernel & \path{65ac095f5cdf4c2e319365c8b0e024d031b19d9fc2b8102e59997afa1e8f9129} \\
P7-T03 operational device suite & \path{d6c818ee29f1a7e659e2f454aec21431d680b3d2d4df048fcf36f4aba87ba22a} \\
P7-T04 propagation profile & \path{87014253023cdb8945ed67f606355d762486884ce3a6de4fa2d32e2af32e2b43} \\
P7-T05 source coupling map & \path{5a9a8f5542a7c8b714bbff7ec06c06449b0c66c0196266051562caf9ce602c6b} \\
P7-T06 variational object & \path{386769e40167c35604625ef7250c027dc1712c82db4790e2895e0b31ac3cfbf7} \\
\bottomrule
\end{tabular}
\end{center}

The eight adopted meanings are, in order: physical source-matter structure;
finite source-matter dynamics; finite operational meanings; source-side
propagation and accessibility; equipped source-matter coupling in its declared
domain; a finite source-matter action; a source-side stress-energy input or
precursor; and componentwise source-side conservation. Each is
\emph{adopted, not derived}.

\subsection{What Position A and P7 do not supply}

The conditional input supplies none of the following:
\begin{enumerate}
  \item realistic continuum matter or a controlled continuum/refinement limit;
  \item Standard Model fields, sectors, representations, or effective sectors;
  \item classical or quantum particle/field statistics;
  \item Hilbert-space, operator-algebraic, or state-algebra structure;
  \item microcausality or target-spacetime locality;
  \item a target tensor $T_{\mu\nu}$ or target matter action;
  \item a derived universal coupling to one source-derived effective geometry;
  \item an equivalence-principle result, $g_{\mathrm{eff}}$, an Einstein-Hilbert
        action, Einstein equations, or a benchmark result.
\end{enumerate}
Realistic matter emergence is explicitly outside the current cycle.

\section{Typed conditional-input contract}

Let $I_{P7}$ be the interface containing only the exact finite types and maps
named above. Let
\[
  \mathcal T_{\rm target}=\{g_{\rm eff},T_{\mu\nu},
  S_{\rm matter}^{\rm target},G_{\mu\nu}=8\pi T_{\mu\nu}\}
\]
denote target-side outputs. A later matter-first packet may read $I_{P7}$ but
must separately construct, type, and justify every map into
$\mathcal T_{\rm target}$.

\paragraph{Conditional-input non-entailment lemma.}
In the free typed interface theory generated by the declared P7 constructors,
if no constructor has codomain in $\mathcal T_{\rm target}$, then no well-typed
term built only from those constructors has codomain in
$\mathcal T_{\rm target}$.

\paragraph{Proof.}
Proceed by structural induction on terms. Variables and constants have only
declared P7 source-side types. Composition, product, restriction, finite sum,
and the declared variations preserve the codomains supplied by their
constituent constructors. Since the generating signature contains no target
codomain and no source-to-target constructor, induction cannot produce a term
of target type. Therefore a target object requires a new typed bridge and its
evidence; a familiar name, protected meaning, validator result, or human
selection cannot supply it. This is a contract-relative type obstruction, not
a global impossibility theorem about future source extensions. $\square$

\section{Evidence-only Gate B and Gate D policy}

For a candidate $c$, let $e$ be its immutable evidence record. Define the Gate
B predicate
\[
Q_B(c,e)=\bigwedge_{i=1}^{8} b_i(c,e)\;\wedge\;r_B(c,e),
\]
where the eight $b_i$ are: target-atlas-free construction; a gauge- and
constraint-reduced hyperbolic principal polynomial; one stable Lorentzian cone
or a physically derived equivalence class; declared-sector universality;
operational clock/rod calibration and scale; tensorial transformation and
local-to-global gluing; finite-variation and coarse-graining robustness; and
absence of a hidden metric-equivalent primitive. The term $r_B$ collects the
required regularity, time-orientation, nondegeneracy, conditioning, and
second-clock-effect analyses.

Define $Q_D(c,e)$ to require a positive Gate B record, complete tensorial
dynamics, universal coupling in scope, conservation or Ward identities,
stable physical degrees of freedom, controlled corrections, and no unresolved
critical independent-review objection. A postulated Einstein-Hilbert action
or postulated Einstein equations cannot satisfy those evidence predicates.

For $G\in\{B,D\}$, the valid evaluator is
\[
\Gate_G(c,e)=
\begin{cases}
\Positive, & Q_G(c,e)\text{ is established by admissible evidence},\\
\Insufficient, & \text{admissible evidence establishes a disqualifying failure},\\
\NotReady, & \text{otherwise}.
\end{cases}
\]

Let $h$ be a protected human decision record. A permissible selector may choose
only from the evidence-supported set $\Supp(e)$. On the extended record define
the evidence projection $\pi_E(c,e,h)=(c,e)$ and require
\[
  \widehat{\Gate}_G=\Gate_G\circ\pi_E.
\]

\paragraph{Evidence-projection invariance theorem.}
For every evidence record $e$, every candidate $c$, and every protected
decision $h$ that does not add independently admissible scientific evidence,
\[
  \widehat{\Gate}_G(c,e,h)=\Gate_G(c,e),\qquad G\in\{B,D\}.
\]
In particular, if any required predicate is absent, protected selection cannot
make the gate positive.

\paragraph{Proof.}
By the factorization requirement, the extended evaluator discards $h$ before
evaluation. Thus both sides evaluate the same pair $(c,e)$. Missing predicates
leave $Q_G$ false or unestablished. Any implementation whose output changes
solely with $h$ fails the factorization equation and is therefore an invalid
Gate B/D evaluator. $\square$

Allowed human authority may choose among evidence-supported candidates, order
repairs, select a scoped evidence-grounded verdict, or decline promotion.
Forbidden authority may not alter an evidence bit, waive a required predicate,
convert a postulate or validator PASS into evidence, declare internal review
external, or turn absent source-to-geometry/source-to-Einstein evidence into a
positive derivation gate.

\section{Position B reopening criteria}

Position B is not active. A separate matter-emergence program may be opened
only when all of the following control predicates are recorded:
\begin{enumerate}
  \item a materially distinct, source-side matter-emergence hypothesis and
        explicit scientific rationale;
  \item a separate registered plan, budget, AgentJob route, claim boundary,
        and protected authority only for any genuinely protected adoption;
  \item an explicit separation from Position A and preservation of the P7
        postulate boundary and historical results;
  \item preregistered target objects for degrees of freedom, internal symmetry,
        statistics, Hilbert/algebraic states, locality, interactions,
        continuum/refinement behavior, and effective Standard Model sectors;
  \item source-provenance, falsification, robustness, and continuum tests plus
        a finite family/reopening budget; and
  \item the same no-postulate Gate B/D rule, including evidence-projection
        invariance.
\end{enumerate}
Opening Position B would authorize research on those obligations, not establish
their truth, adopt a new ontology, or displace Position A automatically.

\section{Disposition and downstream boundary}

P1-T03 is complete at draft/control scope when the exact contract, reopening
criteria, policy, fixtures, source registration, claim-language checks, and
ordinary control receipts pass. The Distance-to-GR ledger is unchanged: P7 is
a scoped source-side postulate input, while $g_{\rm eff}$, target matter
coupling, Einstein equations, Gate B, Gate D, and benchmark promotion remain
blocked. P1-T04 may separately lock Gate B as the sole active physics gate; it
is not executed by this packet.

\end{document}
